\chapter{Estado del Arte}\label{chapter:estado_arte}

\section{Criterios de comparación}\label{sec:criterios_comparacion}

Con el fin de contrastar de manera homogénea las plataformas relevadas en este capítulo, se definió un conjunto de dimensiones funcionales que sintetizan las capacidades centrales abordadas en el marco teórico. Estas dimensiones permiten analizar cada plataforma bajo los mismos parámetros y constituyen la base de la comparativa presentada en la sección~\ref{sec:comparativa_plataformas}. Su selección responde a la naturaleza dual del problema que aborda el presente trabajo, que combina el descubrimiento de videojuegos orientado al jugador con el análisis de su recepción orientado al desarrollador. Los criterios considerados son los siguientes:

El dominio y foco clasifica cada plataforma según el dominio que abarca (videojuegos u otro tipo de contenido) y su orientación principal, ya sea hacia el jugador o hacia el desarrollador; esta distinción resulta central, dado que el dominio de los videojuegos presenta particularidades propias para los sistemas de recomendación (véase la sección~\ref{sec:recomendacion_videojuegos}; \parencite{CheuqueEtAl2019, ViljanenEtAl2020}) y que la analítica orientada a desarrolladores constituye un eje diferenciado (véase la sección~\ref{sec:analitica_audiencias}; \parencite{DrachenEtAl2011}). El tipo de sistema de recomendación evalúa si la plataforma incorpora un sistema de recomendación y, en tal caso, qué enfoque adopta: basado en contenido, colaborativo o híbrido (véase la sección~\ref{sec:sistemas_recomendacion}; \parencite{KoEtAl2022, ShahEtAl2016}). El enfoque de perfilado de usuario considera si la plataforma construye perfiles y mediante qué estrategia, ya sea explícita, implícita o híbrida (véase la sección~\ref{sec:perfilado_usuarios}; \parencite{KanojeEtAl2014, EkeEtAl2019}). El tratamiento de las reseñas distingue entre plataformas que únicamente las recolectan y muestran y aquellas que las procesan mediante técnicas de procesamiento de lenguaje natural para extraer información de las opiniones (véase la sección~\ref{sec:procesamiento_resenas}; \parencite{ZhangEtAl2022, MedhatEtAl2014}). El componente social y de seguimiento atiende a las funcionalidades de interacción que registran la relación entre el usuario y los títulos (calificaciones, listas, seguimiento de otros usuarios y actividad de la comunidad), que constituyen señales explícitas de interés (véase la sección~\ref{sec:recomendacion_videojuegos}; \parencite{CheuqueEtAl2019, ViljanenEtAl2020}). La analítica para desarrolladores determina si la plataforma ofrece información sobre la recepción y el público de sus juegos, y de qué tipo, ya sea de mercado, de comunidad o basada en el análisis de reseñas (véase la sección~\ref{sec:analitica_audiencias}; \parencite{DrachenEtAl2011, GuzsvineczSzucs2023}). Por último, el acceso a los datos releva si la plataforma expone su información de forma pública mediante una interfaz de programación de aplicaciones (API), aspecto relevante en tanto la disponibilidad de datos externos condiciona la construcción de sistemas de recomendación (véase la sección~\ref{sec:recomendacion_videojuegos}; \parencite{ViljanenEtAl2020}).

A partir de estas dimensiones se analiza cada plataforma en las secciones siguientes y se sintetizan los resultados en la comparativa de la sección~\ref{sec:comparativa_plataformas}, lo que permite identificar las oportunidades detalladas en la sección~\ref{sec:oportunidades}.

\section{Plataformas de distribución y recomendación de videojuegos}\label{sec:plataformas_distribucion}

\subsection{Steam}\label{subsec:steam}

Steam, desarrollada por Valve Corporation, es la principal plataforma de distribución digital de videojuegos para PC y funciona simultáneamente como tienda, biblioteca y comunidad. Su catálogo, compuesto por decenas de miles de títulos, constituye un caso paradigmático del problema de sobrecarga de información que motiva a los sistemas de recomendación en este dominio. Por su escala y por la riqueza de los datos de interacción que genera, Steam es el referente obligado para cualquier sistema de recomendación de videojuegos: \textcite{CheuqueEtAl2019} fueron los primeros en estudiar sistemas de recomendación en el contexto de plataformas de videojuegos en línea, tomando a Steam como caso.

En cuanto al sistema de recomendación, Steam ofrece el \emph{Interactive Recommender}, una herramienta que emplea un modelo de aprendizaje automático entrenado sobre los historiales de tiempo de juego de millones de usuarios y que no se apoya directamente en las etiquetas ni en las reseñas, sino en lo que los usuarios efectivamente juegan \parencite{ValveCorporation2020}. Su lógica es de tipo colaborativo: si otros jugadores con hábitos de juego similares a los del usuario disfrutan de un título que este no probó, es probable que también le guste. El usuario puede ajustar los resultados mediante filtros por etiquetas y balancear entre títulos populares o de nicho y entre lanzamientos recientes o clásicos. Steam complementa esta herramienta con la \emph{Discovery Queue}, orientada al contenido y capaz de presentar lanzamientos nuevos que todavía nadie jugó, algo que el \emph{Interactive Recommender} no puede hacer. En el plano académico, \textcite{CheuqueEtAl2019} evaluaron sobre datos de Steam modelos basados en máquinas de factorización y redes neuronales profundas, y hallaron que el modelo de red neuronal profunda obtuvo el mejor desempeño y que el sentimiento extraído de las reseñas resultó menos relevante para la recomendación de lo que cabría esperar.

El perfilado sobre el que opera esta recomendación es predominantemente implícito: se construye a partir de señales de comportamiento como el tiempo de juego, la propiedad de títulos y el historial de compras, más que de preferencias declaradas de forma explícita (véase la sección~\ref{sec:perfilado_usuarios}; \parencite{CheuqueEtAl2019, ViljanenEtAl2020}).

Respecto del tratamiento de las reseñas, Steam permite que los usuarios con tiempo de juego registrado escriban reseñas indicando si recomiendan o no un título, y agrega esas valoraciones en un puntaje de reseñas que se muestra en la página del producto \parencite{ValveCorporationReviews}. Ese puntaje se expresa en categorías que van de \enquote{Extremadamente negativas} a \enquote{Extremadamente positivas}, se calcula esencialmente como el porcentaje de reseñas positivas y distingue entre las reseñas recientes de los últimos 30 días y las de todo el período. El sistema incorpora además salvaguardas frente a campañas coordinadas de reseñas negativas (\emph{review bombing}). No obstante, el procesamiento que realiza Steam se limita a la agregación de la polaridad: la plataforma recolecta, muestra y resume el signo de las reseñas, pero no analiza su contenido a nivel de aspectos ni extrae de forma automática las opiniones sobre características específicas del juego (véase la sección~\ref{sec:procesamiento_resenas}).

En cuanto al componente social y de seguimiento, Steam dispone de listas de deseos, etiquetas asignadas por la comunidad, perfiles, foros y funciones sociales, que enriquecen tanto el perfil del usuario como la caracterización de los títulos. Desde la perspectiva del desarrollador, expone el puntaje de reseñas y métricas asociadas al desempeño del título, pero no un análisis de las opiniones desagregado por aspecto ni una caracterización de la recepción más allá de la polaridad general; ese tipo de análisis queda cubierto por herramientas de terceros que se analizan en la sección~\ref{sec:herramientas_analitica}. Finalmente, en lo relativo al acceso a los datos, Steam expone parte de su información mediante la Steam Web API, lo que habilita a terceros (incluido el presente trabajo) a construir aplicaciones sobre sus datos \parencite{ValveCorporationReviews}.

En síntesis, Steam combina distribución, recomendación y reseñas en una sola plataforma, pero su recomendación se apoya en señales implícitas de comportamiento y deja de lado tanto el contenido textual de los títulos como el perfilado explícito, mientras que su tratamiento de las reseñas no supera la agregación de polaridad. Estas limitaciones delimitan parte del espacio de oportunidad del presente trabajo, que incorpora filtrado basado en contenido y un análisis de reseñas orientado al desarrollador.

\section{Plataformas de registro, reseñas y comunidad}\label{sec:plataformas_registro}

\subsection{Backloggd}\label{subsec:backloggd}

Backloggd es una plataforma web de registro y seguimiento de videojuegos que funciona, en esencia, como un diario o bitácora de juego: permite a los usuarios registrar, puntuar y reseñar los juegos que jugaron, y agregar a una lista de pendientes (\emph{backlog}) aquellos que desean jugar \parencite{Backloggd}. La propia plataforma se presenta de forma explícita como el equivalente para videojuegos de los servicios de catalogación de medios, en la línea de Goodreads para libros o Letterboxd para películas \parencite{Backloggd}.

En cuanto al perfilado y los datos de interacción, Backloggd se apoya en señales explícitas declaradas por el propio usuario: este marca cada título como \enquote{jugando}, \enquote{completado} o \enquote{en backlog}, lo puntúa y escribe reseñas, construyendo así un historial de su actividad de juego (véase la sección~\ref{sec:perfilado_usuarios}; \parencite{KanojeEtAl2014}). A diferencia de Steam, no infiere el perfil a partir del comportamiento registrado de manera automática, sino que depende de la carga manual del usuario, lo que lo ubica como un caso de perfilado predominantemente explícito.

El componente social y de seguimiento constituye el núcleo de la plataforma: permite seguir a otros jugadores, leer y comparar sus reseñas y calificaciones, y descubrir juegos a través de listas y \emph{feeds} de actividad. Las listas son de formato libre, lo que habilita recopilaciones temáticas creadas por la comunidad.

Respecto del tratamiento de las reseñas, Backloggd se limita (al igual que Steam) a recolectar, mostrar y permitir comparar reseñas y puntuaciones, sin realizar ningún análisis automático de su contenido (véase la sección~\ref{sec:procesamiento_resenas}). Tampoco incorpora un sistema de recomendación algorítmico: el descubrimiento de títulos es social y manual, mediado por las listas, los perfiles seguidos y la actividad de la comunidad, y no por un modelo que prediga afinidad. Por su orientación, carece asimismo de todo módulo de analítica dirigido a desarrolladores. Finalmente, en cuanto al acceso a los datos, su catálogo de juegos proviene de una base de datos de videojuegos de terceros y la plataforma no ofrece una API pública oficial para la extracción de datos.

En síntesis, Backloggd aporta un modelo sólido de catalogación social, seguimiento y construcción de un perfil explícito del jugador (dimensiones poco desarrolladas en Steam), pero no ofrece recomendación personalizada, análisis de reseñas ni información para desarrolladores. Su valor como referencia para el presente trabajo radica en ese enfoque de perfilado explícito y comunidad, que complementa el modelo implícito y algorítmico de las plataformas de distribución.

\subsection{RAWG}\label{subsec:rawg}

RAWG es un servicio de descubrimiento de videojuegos construido sobre la que se presenta como la mayor base de datos abierta del rubro: reúne más de 350.000 títulos en más de 50 plataformas, con metadatos ricos como etiquetas, géneros, desarrolladores, editores, fechas de lanzamiento, puntajes de Metacritic, enlaces a tiendas y tiempo medio de juego \parencite{RAWG}. Sobre esa base ofrece, al igual que las plataformas anteriores, funciones de registro y comunidad: permite reunir los juegos en un perfil, calificarlos, reseñarlos y ver qué están jugando otros usuarios \parencite{RAWG}.

Su rasgo más distintivo en cuanto al sistema de recomendación es que incorpora un algoritmo de aprendizaje automático que analiza el contenido del propio juego para sugerir títulos similares. A diferencia del enfoque colaborativo basado en comportamiento de Steam, RAWG aplica así una recomendación basada en contenido (véase la sección~\ref{sec:sistemas_recomendacion}; \parencite{KoEtAl2022}), que es precisamente uno de los enfoques que adopta el presente trabajo.

En cuanto al tratamiento de las reseñas, RAWG recolecta calificaciones y reseñas de los usuarios pero (nuevamente como Steam y Backloggd) no realiza un análisis automático de su contenido (véase la sección~\ref{sec:procesamiento_resenas}). Su componente social y de seguimiento es comparable al de Backloggd, con perfiles, colecciones y seguimiento de otros jugadores, aunque subordinado a su naturaleza de base de datos.

La dimensión más diferencial de RAWG es el acceso a los datos. La plataforma expone una API REST pública con autenticación por clave, que permite buscar y filtrar juegos por plataforma, género, desarrollador, editor, etiqueta, fecha y calificación, y es de uso gratuito para proyectos personales y usos comerciales pequeños, con atribución obligatoria \parencite{RAWG}. Esto la convierte en una fuente de datos de referencia para terceros que construyen aplicaciones de descubrimiento y motores de recomendación, función que en el ecosistema cumple junto a la Steam Web API utilizada en el presente trabajo. En contraste, RAWG no ofrece un módulo de analítica orientado a desarrolladores: a estos les sirve como proveedor de datos, no como herramienta de análisis de la recepción de sus títulos.

En síntesis, RAWG combina una base de datos amplia y abierta, recomendación basada en contenido y una API pública robusta, lo que lo acerca al enfoque de contenido del presente trabajo y lo distingue de Backloggd por su accesibilidad de datos. No obstante, comparte con las plataformas anteriores la ausencia de análisis de reseñas y de información sobre la recepción dirigida a desarrolladores.

\section{Plataforma de referencia de otro dominio}\label{sec:plataforma_referencia}

\subsection{Letterboxd}\label{subsec:letterboxd}

Letterboxd no pertenece al dominio de los videojuegos, sino al del cine, pero se incluye en este relevamiento por constituir la principal referencia e inspiración de las plataformas de catalogación social de medios; de hecho, varias plataformas del rubro de los videojuegos, Backloggd entre ellas, se presentan explícitamente como su equivalente. Se define como una red social de descubrimiento cinematográfico que funciona como un diario para registrar y compartir la opinión sobre las películas a medida que se ven, y que permite calificarlas, reseñarlas y etiquetarlas, seguir a otros usuarios, mantener una lista de pendientes y crear listas sobre cualquier tema \parencite{Letterboxd}. Las calificaciones utilizan una escala de media a cinco estrellas, y la plataforma se ha consolidado como un fenómeno cultural ampliamente adoptado dentro de su dominio.

En términos de perfilado, Letterboxd es un caso paradigmático de perfil explícito: el usuario declara activamente sus gustos mediante calificaciones, reseñas, diario y listas, lo que lo aproxima al modelo de Backloggd más que al de las plataformas de distribución (véase la sección~\ref{sec:perfilado_usuarios}; \parencite{KanojeEtAl2014}). Ese perfil rico, junto con su componente social y de seguimiento, constituye su mayor fortaleza y el principal motor de descubrimiento de la plataforma.

Ahora bien, ese descubrimiento se apoya fundamentalmente en mecanismos sociales y editoriales, la actividad de los usuarios seguidos, las listas de la comunidad, las listas destacadas curadas por género, estudio o temática y las calificaciones agregadas, más que en un sistema de recomendación personalizada basado en algoritmos. En cuanto al tratamiento de las reseñas, y pese a que estas influyen de manera relevante en las decisiones de consumo de los usuarios (véase la sección~\ref{sec:procesamiento_resenas}; \parencite{AliEtAl2016}), Letterboxd las recolecta, muestra y permite compartir, pero no analiza su contenido de forma automática. Tampoco ofrece analítica orientada a los creadores de las obras.

Es decir, Letterboxd representa el estándar de excelencia en catalogación social, perfilado explícito y comunidad para un dominio de medios, y es por ello la referencia de diseño del presente trabajo. Sin embargo, incluso esta plataforma de referencia deja sin cubrir dos capacidades centrales para los objetivos planteados: la recomendación personalizada mediante algoritmos y el análisis automático de las reseñas, esta última especialmente relevante para un módulo orientado a desarrolladores. El presente trabajo retoma de Letterboxd el modelo de experiencia y de perfilado, y lo extiende con esas dos capacidades en el dominio de los videojuegos.

\section{Herramientas de analítica para desarrolladores}\label{sec:herramientas_analitica}

A diferencia de las plataformas anteriores, orientadas principalmente al jugador, existe un ecosistema de herramientas dirigidas a los desarrolladores y editores cuyo propósito es transformar los datos generados en torno a los videojuegos en información útil para la toma de decisiones (véase la sección~\ref{sec:analitica_audiencias}; \parencite{DrachenEtAl2011}). Dentro de este ecosistema pueden distinguirse dos grandes grupos según el tipo de dato que explotan: las herramientas de analítica de mercado y las de análisis de reseñas.

\subsection{Analítica de mercado}\label{subsec:analitica_mercado}

Un primer conjunto de herramientas se especializa en estimar y monitorear el desempeño comercial de los títulos. Servicios como SteamSpy, SteamDB, Gamalytic y Video Game Insights ofrecen datos sobre la cantidad de propietarios, los ingresos estimados, los jugadores concurrentes, las listas de deseos y la popularidad por etiquetas, a partir de los datos públicos de Steam \parencite{SteamSpy, VideoGameInsights}. Este tipo de analítica responde a la dimensión de mercado y telemetría descrita en el marco teórico, donde los datos de cliente y de comportamiento permiten orientar decisiones de negocio (véase la sección~\ref{sec:analitica_audiencias}; \parencite{DrachenEtAl2011}).

Cabe aclarar que esta clase de herramientas opera sobre agregados de ventas y comportamiento, no sobre el contenido de las opiniones de los jugadores. El presente trabajo no aborda la analítica de mercado: se la incluye únicamente para delimitar el alcance del ecosistema y diferenciarla del enfoque adoptado, centrado en el análisis de las reseñas.

\subsection{Análisis de reseñas}\label{subsec:analisis_resenas}

Un segundo conjunto de herramientas y enfoques se ocupa específicamente de las reseñas de los jugadores como fuente de información para los desarrolladores. Como se señaló (véase la sección~\ref{subsec:steam}), tanto Steam como la mayoría de las plataformas reducen las reseñas a una clasificación binaria (recomendado o no recomendado) y a un puntaje agregado de polaridad, lo que deja de lado los sentimientos matizados que los jugadores expresan sobre características específicas del juego. Frente a esta limitación, existen herramientas que analizan las reseñas con mayor profundidad: algunas las filtran y resumen por región o idioma, como SteamData.ninja o Steam Review Explorer, y otras aplican técnicas de procesamiento de lenguaje natural para clasificar el sentimiento de los textos.

El nivel más granular de este análisis es el análisis de sentimiento basado en aspectos (ABSA), que identifica las opiniones sobre características concretas (jugabilidad, gráficos, historia o música) en lugar de un único sentimiento global (véase la sección~\ref{subsec:absa}; \parencite{ZhangEtAl2022}). Diversos trabajos han mostrado el valor de este enfoque aplicado a reseñas de Steam, por ejemplo al identificar qué aspectos resuenan más entre los jugadores, como la narrativa, para orientar el desarrollo, o al caracterizar la relación entre la extensión y el sentimiento de las reseñas según el género \parencite{GuzsvineczSzucs2023}. Sin embargo, este tipo de análisis a nivel de opinión se encuentra mayormente en trabajos académicos y proyectos de investigación, más que en productos consolidados orientados a desarrolladores, donde predomina la analítica de mercado o, a lo sumo, el resumen agregado del sentimiento.

Es precisamente en este espacio donde se posiciona el módulo orientado a desarrolladores del presente trabajo: aplicar el análisis de las reseñas para extraer información estructurada sobre la recepción de cada título, más allá de la polaridad agregada, acercando a los desarrolladores una comprensión desagregada de las opiniones de sus jugadores.

\section{Comparativa entre plataformas}\label{sec:comparativa_plataformas}

A partir del análisis de las secciones anteriores, y tomando como referencia las dimensiones definidas en la sección~\ref{sec:criterios_comparacion}, es posible contrastar las plataformas relevadas y reconocer patrones comunes.

En cuanto al sistema de recomendación, solo Steam y RAWG incorporan mecanismos algorítmicos, y lo hacen mediante enfoques opuestos: Steam se apoya en un modelo colaborativo basado en el comportamiento de los usuarios, mientras que RAWG aplica un enfoque basado en el contenido de los juegos. Backloggd y Letterboxd, en cambio, no ofrecen recomendación algorítmica y delegan el descubrimiento en mecanismos sociales y editoriales. Ninguna de las plataformas combina ambos enfoques de manera híbrida y personalizada, tal como se propone en el presente trabajo.

En el perfilado de usuario se observa una división clara según el tipo de plataforma. Las orientadas a la distribución, como Steam, construyen el perfil a partir de señales implícitas derivadas del comportamiento, mientras que las de catalogación social, como Backloggd y Letterboxd, se basan en perfiles explícitos declarados por el propio usuario mediante calificaciones, reseñas y listas. Sin embargo, ninguna integra ambos tipos de señal en una representación combinada del usuario.

El patrón más significativo aparece en el tratamiento de las reseñas. Todas las plataformas analizadas las recolectan y las muestran (y, en el caso de Steam, las resumen en un puntaje de polaridad), pero ninguna analiza su contenido a nivel de opinión. Ese análisis solo se encuentra en herramientas y proyectos orientados a desarrolladores, y allí mayormente en el ámbito de la investigación, desacoplado tanto del entorno de recomendación como de la experiencia de cara al jugador. De manera complementaria, la analítica para desarrolladores existe como un segmento separado, dominado por la información de mercado, en el que el análisis de la opinión de los jugadores permanece desatendido como producto consolidado.

El componente social y de seguimiento, por su parte, está presente en todas las plataformas orientadas al jugador, y constituye la fortaleza principal de Backloggd y Letterboxd, pero se encuentra ausente en las herramientas dirigidas a desarrolladores. Finalmente, en cuanto al acceso a los datos, Steam y RAWG exponen interfaces de programación públicas que habilitan su uso por terceros, mientras que Backloggd no ofrece una API pública oficial y la de Letterboxd es de acceso restringido.

La tabla~\ref{tab:comparativa_plataformas} sintetiza el análisis desarrollado en las secciones anteriores, contrastando cada plataforma relevada según las dimensiones definidas en la sección~\ref{sec:criterios_comparacion} e incorporando, en la última fila, la propuesta del presente trabajo.

\setlength{\tabcolsep}{4pt}
\begin{longtable}{|>{\raggedright\arraybackslash}p{1.8cm}|>{\raggedright\arraybackslash}p{2.3cm}|>{\raggedright\arraybackslash}p{1.5cm}|>{\raggedright\arraybackslash}p{2.5cm}|>{\raggedright\arraybackslash}p{2.5cm}|>{\raggedright\arraybackslash}p{1.9cm}|}
	\caption{Comparación de las plataformas relevadas.}\label{tab:comparativa_plataformas} \\
	\hline
	\textbf{Plataforma} & \textbf{Recomendación} & \textbf{Perfilado} & \textbf{Tratamiento de reseñas} & \textbf{Analítica para desarrolladores} & \textbf{Acceso a los datos} \\
	\hline
	\endfirsthead
	\multicolumn{6}{l}{\textit{Continuación de la Tabla \ref{tab:comparativa_plataformas}}} \\
	\hline
	\textbf{Plataforma} & \textbf{Recomendación} & \textbf{Perfilado} & \textbf{Tratamiento de reseñas} & \textbf{Analítica para desarrolladores} & \textbf{Acceso a los datos} \\
	\hline
	\endhead
	\hline
	\endfoot
	\hline
	\endlastfoot
	Steam & Colaborativo, basado en tiempo de juego & Implícito & Agregación de polaridad & Métricas de desempeño del título & API pública \\
	Backloggd & No algorítmica; social & Explícito & Solo recolecta y muestra & No ofrece & Sin API pública \\
	RAWG & Basada en contenido & Explícito & Solo recolecta y muestra & No ofrece & API REST pública \\
	Letterboxd (cine) & No algorítmica; social y editorial & Explícito & Solo recolecta y muestra & No ofrece & Acceso restringido \\
	SteamSpy, SteamDB, Gamalytic, VG Insights & No aplica & No aplica & No procesa el contenido & Analítica de mercado & API pública \\
	GameTrack & Híbrida: contenido, colaborativo y respaldo por popularidad & Explícito e implícito & Análisis de sentimiento basado en aspectos & Recepción desagregada por aspecto & No expone datos a terceros \\
\end{longtable}
\setlength{\tabcolsep}{6pt}

La lectura de la tabla permite reconocer que ninguna plataforma reúne simultáneamente las dos capacidades centrales del presente trabajo. Las que incorporan recomendación algorítmica lo hacen mediante un único enfoque, ya sea colaborativo o basado en contenido, y ninguna combina ambos ni integra señales explícitas e implícitas en un mismo perfil. En cuanto al tratamiento de las reseñas, todas se limitan a recolectarlas y mostrarlas, y a lo sumo a resumir su polaridad, sin analizar el contenido de las opiniones. Por último, la analítica orientada al desarrollador existe únicamente en herramientas separadas del entorno de cara al jugador y se restringe a información de mercado. El valor diferencial de la propuesta no reside entonces en una capacidad inédita, sino en la integración de capacidades que actualmente se encuentran distribuidas entre productos distintos.

En conjunto, ninguna de las plataformas relevadas reúne en un mismo producto una recomendación personalizada para el jugador y un análisis de las reseñas orientado al desarrollador. Esa ausencia delimita el espacio de oportunidad que se detalla en la sección siguiente.

\section{Oportunidades detectadas}\label{sec:oportunidades}

El análisis comparativo de las secciones anteriores permite identificar un conjunto de oportunidades que fundamentan el desarrollo del presente trabajo y que se alinean con sus objetivos: ofrecer recomendaciones personalizadas al jugador y brindar a los desarrolladores información sobre la recepción de sus títulos.

La primera oportunidad se encuentra en la recomendación. Como se observó, las plataformas que incorporan recomendación algorítmica adoptan un único enfoque (Steam, colaborativo basado en comportamiento; RAWG, basado en contenido), sin combinarlos. Esto deja espacio para un sistema híbrido que integre el filtrado basado en contenido y el filtrado colaborativo, aprovechando las ventajas de cada uno y mitigando sus limitaciones, en particular el problema del arranque en frío (véase las secciones~\ref{subsec:enfoque_hibrido} y~\ref{subsec:aplicacion_perfilado}). La incorporación de un mecanismo de respaldo basado en popularidad para los casos de información insuficiente refuerza esta posibilidad, permitiendo generar recomendaciones útiles incluso en las primeras etapas de uso.

Una segunda oportunidad surge en el aprovechamiento del perfil del usuario. Mientras que las plataformas de distribución se apoyan en señales implícitas y las de catalogación social en señales explícitas, ninguna integra ambos tipos de información. La construcción de perfiles dentro de la propia plataforma, a partir de las interacciones del usuario (calificaciones, reseñas e historial), habilita una representación más completa que sirva de base a la recomendación personalizada.

La tercera oportunidad, y la más diferencial, se ubica en el análisis de las reseñas orientado al desarrollador. El relevamiento mostró que, si bien todas las plataformas recolectan reseñas, ninguna analiza su contenido a nivel de opinión, y que el análisis de la recepción a partir de las reseñas se encuentra desatendido como producto consolidado, separado del entorno de cara al jugador. Esto abre la posibilidad de un módulo que procese las reseñas para extraer información estructurada sobre cómo es recibido cada título y qué tipo de jugadores lo valoran, aportando a los desarrolladores una comprensión desagregada que excede la polaridad agregada (véase las secciones~\ref{sec:procesamiento_resenas} y~\ref{sec:analitica_audiencias}).

Finalmente, la oportunidad transversal que articula a las anteriores consiste en reunir estas capacidades en un mismo producto. Ninguna de las plataformas relevadas combina una recomendación personalizada de cara al jugador con un análisis de las reseñas orientado al desarrollador; tomando como referencia la experiencia de catalogación social de plataformas como Letterboxd y Backloggd, y extendiéndola con la recomendación híbrida y el análisis de reseñas ausentes en ellas, el presente trabajo se posiciona en ese espacio no cubierto, atendiendo de forma conjunta tanto a los jugadores como a los desarrolladores.
